\SourceLeaf{286}%
one does need approximate values for the lengths of the lines, the oriented azimuths, and the latitudes of the endpoints. These can be obtained only by a preliminary calculation of the triangles. This is of no consequence, however, since it is in many respects advisable anyway to carry out preliminary calculations step by step as the measurements proceed. Such calculations are even necessary for centering the angles measured off center and for determining the spherical or spheroidal excess of the sum of the angles of each triangle. Indeed, for the first purpose, a very rough approximation will always suffice, given how small those reductions are, whereas precise centering may sometimes require a much closer approximation if the stations are appreciably off center. I have developed the rules so that, \textit{if} one wants to take these reductions into account, everything needed to calculate them as accurately as possible is at hand, and if one does \textit{not} want to take them into account, one can readily see exactly how little is lost by leaving them out. Throughout the Hanoverian triangulation network, the reductions are so extremely small that taking them into account appears entirely unnecessary. Across the full twelve degrees of latitude for which I supply the auxiliary material, they still fall below the fractions of a second of arc to which one ordinarily limits the calculations for most measurements. To make this especially clear, I include the numerical calculations for a couple of examples here.

In the Hanoverian triangulation network, the greatest reductions occur in the directions of the sides of the triangle Brocken-Hohehagen-Inselsberg. This is both the largest triangle and the farthest from the standard parallel. For all the other triangle sides, the reductions nowhere exceed two thousandths of a second, and most do not even reach \(0''001\).

For these points we have
\begin{center}
\begin{tabular}{l|c|c|r}
&\multicolumn{2}{c|}{Latitude}&\\
&on the ellipsoid&on the sphere&\(k\)\\\hline
Brocken&\(51^0\quad48'\quad2''\)&\(51^0\quad46'\quad3''\)&\(0''164\)\\
Hohehagen&\(51\quad28\quad31\)&\(51\quad26\quad35\)&\(0{,}303\)\\
Inselsberg&\(50\quad51\quad9\)&\(50\quad49\quad16\)&\(0{,}687\)
\end{tabular}
\end{center}
